\documentclass[12pt,twoside,a4paper]{article}

\usepackage{amsmath,amssymb}
\usepackage[margin=2.5cm]{geometry}
\usepackage{relsize}
\usepackage{multicol}

\usepackage{color}
\usepackage[most]{tcolorbox}

% This package is only included to let me comment the other exercises in the document, so I can preview the one that I am actively working on as if it were the first, so I don't have to scroll down to see the result.
% I should remove it when I finish the document
\usepackage{comment}

% This one lets me use labels in the enumerate environment, so I can reference them later in the document
\usepackage{enumitem}

\title{
  \textbf{Linear Algebra and Analytic Geometry}\\[0.35em]
  \large Unit 1: Matrices and Systems\\[0.2em]
  \normalsize Worked Solutions to Exercises 1--34
}

\author{Alex Vergara}
%\date{\today}

\newenvironment{banswer}
  {\color{black}}
  {}
\newenvironment{answer}
  {\color{blue}}
  {}
\newenvironment{ranswer}
  {\color{red}}
  {}

\begin{document}

\maketitle

\begin{enumerate}

% ---------- Exercise 1 --------------
\item Consider the following matrix: $A = \begin{bmatrix}
11 & \frac{1}{2} &  -2 & 3 \\
4 & 1 & 1 & 0
\end{bmatrix}.
$

    \begin{enumerate}
        \item What type of matrix is $A$?
                \begin{answer}
                    Is a 2 by 4 matrix $\left(2 \times 4\right)$, so it's a rectangular matrix
                \end{answer}
        \item What are the entries of $A$ at positions $\left(1, 2\right)$ and $\left(2, 1\right)$?
                \begin{answer}
                    \begin{itemize}
                        \item $\mathlarger{a}_{1,2} = \frac{1}{2}$
                        \item $\mathlarger{a}_{2,1} = 4$
                    \end{itemize}
                \end{answer}
        \item At which positions in $A$ is the entry equal to 1?
                \begin{answer}
                    In $a_{2,2}$ and $a_{2,3}$
                \end{answer}
    \end{enumerate}

% ---------- Exercise 2 --------------
\item  Write the matrix $\left(A=\left[a_{ij}\right]_{3\times3}\right)$ defined by:
        \begin {enumerate}
            \item $\mathlarger{a}_{ij} = i + j$
            \begin{answer}
                $$
                A = \begin{pmatrix}
                2 & 3 & 4 \\
                3 & 4 & 5 \\
                4 & 5 & 6
                \end{pmatrix}
                $$
            \end{answer}
            \item $\mathlarger{a}_{ij} = \begin{cases}
                    0, & \text{if } i = j,\\
                    -1, & \text{if } i > j,\\
                    2, & \text{if } i < j.
                    \end{cases}$
            \begin{answer}
                $$
                A = \begin{pmatrix}
                0 & 2 & 2 \\
                -1 & 0 & 2 \\
                -1 & -1 & 0
                \end{pmatrix}
                $$
            \end{answer}
        \end{enumerate}


% ---------- Exercise 3 --------------
\item For each of the following parts, determine, if possible, the real numbers $x$ and $y$ such that the matrices $A$ and $B$ are equal.

    \textcolor{red}{Remember}: Two matrices are equal if they have the \underline{same dimensions} and all their corresponding entries \textbf{are equal}.

    \begin{enumerate}
         % ------------------------ (a) ------------------------
        \item
            $A = \begin{bmatrix}
                1 & 0 \\
                x & y \\
                -1 & 3
            \end{bmatrix}$, \quad
            $B = \begin{bmatrix}
                1 & 0 \\    
                \frac{y}{2} & x + 1 \\    
                -1 & 3 \\    
            \end{bmatrix}$
    
    \textbf{Solution:}
    \begin{banswer}
        \begin{itemize}
            \item From the first row: $\mathlarger{a}_{11} = 1 = \mathlarger{b}_{11}$ and $\mathlarger{a}_{12} = 0 = \mathlarger{b}_{12}$.
            \item From the second row: $\mathlarger{a}_{21} = x = \frac{y}{2} = \mathlarger{b}_{21}$ and $\mathlarger{a}_{22} = y = x + 1 = \mathlarger{b}_{22}$.
            \item From the third row: $\mathlarger{a}_{31} = -1 = \mathlarger{b}_{31}$ and $\mathlarger{a}_{32} = 3 = \mathlarger{b}_{32}$.
        \end{itemize}

        The point is that rows 1 and 3 do not give us any information about $x$ and $y$, so we only need to solve the system of equations from the second row:
        $$
        \left\{
        \begin{aligned}
            x &= \frac{y}{2},\\
            y &= x+1.
        \end{aligned}
        \right.
        $$

        Substitute $x=\frac{y}{2}$ into the second equation:
        $$
        \begin{aligned}
            y &= \frac{y}{2}+1,\\
            2y &= y+2,\\
            y &= 2,\qquad x=\frac{y}{2}=1.
        \end{aligned}
        $$

        \begin{tcolorbox}
            \centering
            $x=1,\qquad y=2$
        \end{tcolorbox}
    \end{banswer}

% ------------------------ (b) ------------------------
    \item 
    $A = \begin{bmatrix}
        x^3 & 2 & 0 \\
        3 & y & -1
    \end{bmatrix}$, \quad
    $B = \begin{bmatrix}
        3x + 1 & 2 & y \\
        3 & -2 & -1
    \end{bmatrix}$

    \textbf{Solution:}
    \begin{banswer}
    \end{banswer}

% ------------------------ (c) ------------------------
    \item 
    $A = \begin{bmatrix}
        2 & -5 \\
        0 & x + y \\
        2y & 3
    \end{bmatrix}$, \quad
    $B = \begin{bmatrix}
        2 & -5 \\    
        0 & 2x + 4 \\    
        y + 1 & 1 \\    
    \end{bmatrix}$

    \textbf{Solution:}

    \end{enumerate}

\item Consider the following matrices:\\[1em]
    \begin{array}{ccc}
        A = \begin{bmatrix} 2 & 2 \\ 3 & 3 \\ -1 & -1 \end{bmatrix},
        &
        B = \begin{bmatrix} -1 & 3 \\ -4 & 0 \\ 0 & -1 \end{bmatrix},
        &
        C = \begin{bmatrix} 2 & -1 \\ 3 & 0 \\ \end{bmatrix},
        \\[2em]
        D = \begin{bmatrix} 1 & 2 & 3 \\ 3 & 1 & -2 \\ \end{bmatrix},
        &
        E = \begin{bmatrix} 1 & 0 & 2 \\ -1 & 2 & 1 \\ -2 & 2 & 4 \\ \end{bmatrix}, &
        F = \begin{bmatrix} -2 & 1 & 3 \\ -1 & 1 & -2 \\ 0 & 1 & 3 \\ \end{bmatrix}
    \end{array}
    \\[2em]
    Determine the following matrices, only if they are defined. If they are not defined, explain why.

    \begin{multicols}{4}
        \begin{enumerate}[label=\Alph*)]
            \item $A+B$
            \item $B+A$
            \item $A-B$
            \item $D+E$
            \item $-A+2B$
            \item $A+C$
            \item $A+D$
            \item $A-A$
            \item $C-C$
            \item $AB$
            \item $AC$
            \item $CA$
            \item $AD$
            \item $DA$
            \item $EF$
            \item $FE$
            \item $A(3D)$
            \item $(AC)D$
            \item $A(CD)$
            \item $(2A)(C(3D))$
            \item $(2A)((2C)(2D))$
            \item $A(CD+D)$
            \item $C^3$
        \end{enumerate}
        \end{multicols}
        % Answers
        \begin{enumerate}[label=\Alph*)]
            \begin{multicols}{3}
                \item \begin{answer} % A)
                $\begin{bmatrix}
                    1 & 5 \\
                    -1 & 3 \\
                    -1 & -2
                \end{bmatrix}$
           \end{answer}, \par
           \item \begin{answer} % B)
                $\begin{bmatrix}
                    1 & 5 \\
                    -1 & 3 \\
                    -1 & -2
                \end{bmatrix}$
           \end{answer}
           \item \begin{answer} % C)
                $\begin{bmatrix}
                    3 & -1 \\
                    7 & 3 \\
                    -1 & 0
                \end{bmatrix}$
           \end{answer}
            \end{multicols}
           \item \begin{ranswer} % D)
                Not defined, because $D$ is a $2 \times 3$ matrix and $E$ is a $3 \times 3$ matrix, so they do not have the same dimensions.
           \end{ranswer}
           \item \begin{answer} % E)
                $-1 \space \cdot \begin{bmatrix}
                    2 & 2 \\
                    3 & 3 \\
                    -1 & -1
                \end{bmatrix}$ + $2 \space \cdot \begin{bmatrix}
                    -1 & 3 \\
                    -4 & 0 \\
                    0 & -1
                \end{bmatrix}$ =
                $\begin{bmatrix}
                    -4 & 5 \\
                    -11 & 3 \\
                    0 & -2
                \end{bmatrix}$
           \end{answer}
           \item \begin{ranswer} % F)
                Not defined, because $A_{3 \times 2} = [a_{ij}] = A$ and $C$ is a $C_{2 \times 2}$ matrix, so they do not have the same dimensions.
           \end{ranswer}
           \item \begin{ranswer} % G)
                Not defined, because $A$ is a $3 \times 2$ matrix and $D$ is a $2 \times 3$ matrix, so they do not have the same dimensions.
           \end{ranswer}
           \item \begin{answer} % H)
                $H = 0_{3x3}$
              \end{answer}
            \item \begin{answer} % I)
                $H = 0_{2x2}$
              \end{answer}
        \end{enumerate}
    % ---------- End list of exercises --------------
\end{enumerate}
\end{document}
